\section*{Resumen}
Se desarrolló el diseño de las fundaciones de un edificio de hormigón armado de 21 pisos más un subterráneo, ubicado en Antofagasta, zona sísmica 3, suelo tipo A, a partir de las solicitaciones estáticas y sísmicas del modelo estructural. Se diseñaron 30 elementos: 21 muros resueltos mediante zapatas corridas, verificadas en una dirección, y 9 columnas resueltas mediante zapatas aisladas, verificadas de forma biaxial. Los cinco muros perimetrales del sector, P1, PN, P13, Plb y P12, se resuelven mediante zapata excéntrica, dado que su base no puede desarrollarse hacia el exterior de la línea de cierro del edificio. Todas las fundaciones adoptan una altura uniforme $H = 1,20$~m y se diseñan sobre un suelo con presión admisible de 70~tonf/m$^2$ en condición estática y 93~tonf/m$^2$ en condición sísmica. En P1, PN, P13 y Plb la excentricidad geométrica reduce la fracción de base comprimida a un rango de 67\% a 68\% y eleva la presión de contacto al orden de 33 a 35~tonf/m$^2$, con factores de seguridad cercanos a 2,0; P12, pese a su misma condición de borde, presenta una presión considerablemente menor, de 11,14~tonf/m$^2$, con la base en compresión completa. En los muros interiores más solicitados, como PE y PJ, la presión sísmica resulta determinante y el diseño queda ajustado al límite admisible, con factores de seguridad cercanos a 1,00. La zona comprimida de los elementos con zapata centrada supera el 80\% de la base en todos los casos, con un mínimo de 91,6\% en el muro PH. Las verificaciones de corte y punzonamiento resultan holgadas dada la altura útil de las zapatas, y la estabilidad al deslizamiento se satisface en su totalidad, con los casos más ajustados en los cinco muros perimetrales con zapata excéntrica, apenas por sobre el mínimo de 1,5. La armadura queda gobernada por la cuantía mínima, con una disposición de $\phi 18 @ 100$~mm, disponiéndose una parrilla de $\phi 10 @ 200$~mm en los elementos de mayor solicitación de flexión. Todos los elementos verifican los requisitos de resistencia conforme a ACI 318-19, NCh433 y DS 60.